\documentclass[10pt,a4paper]{amsart}
\usepackage{amsmath,amssymb,amsthm,graphicx}
\usepackage{verbatim}
\newenvironment{reference}{\comment}{\endcomment}
%\newenvironment{reference}{}{}
\newenvironment{slogan}{\comment}{\endcomment}
\newenvironment{history}{\comment}{\endcomment}

% For commutative diagrams we use Xy-pic
\usepackage[all]{xy}

% We use 2cell for 2-commutative diagrams.
\xyoption{2cell}
% Theorem environments.
%
\theoremstyle{plain}
\newtheorem{theorem}[subsection]{Theorem}
\newtheorem{proposition}[subsection]{Proposition}
\newtheorem{lemma}[subsection]{Lemma}

\theoremstyle{definition}
\newtheorem{definition}[subsection]{Definition}
\newtheorem{example}[subsection]{Example}
\newtheorem{exercise}[subsection]{Exercise}
\newtheorem{situation}[subsection]{Situation}

\theoremstyle{remark}
\newtheorem{remark}[subsection]{Remark}
\newtheorem{remarks}[subsection]{Remarks}

\numberwithin{equation}{subsection}

% Macros
%
\def\lim{\mathop{\mathrm{lim}}\nolimits}
\def\colim{\mathop{\mathrm{colim}}\nolimits}
\def\Spec{\mathop{\mathrm{Spec}}}
\def\Hom{\mathop{\mathrm{Hom}}\nolimits}
\def\Ext{\mathop{\mathrm{Ext}}\nolimits}
\def\SheafHom{\mathop{\mathcal{H}\!\mathit{om}}\nolimits}
\def\SheafExt{\mathop{\mathcal{E}\!\mathit{xt}}\nolimits}
\def\Sch{\mathit{Sch}}
\def\Mor{\mathop{\mathrm{Mor}}\nolimits}
\def\Ob{\mathop{\mathrm{Ob}}\nolimits}
\def\Sh{\mathop{\mathit{Sh}}\nolimits}
\def\NL{\mathop{N\!L}\nolimits}
\def\CH{\mathop{\mathrm{CH}}\nolimits}
\def\proetale{{pro\text{-}\acute{e}tale}}
\def\etale{{\acute{e}tale}}
\def\QCoh{\mathit{QCoh}}
\def\Ker{\mathop{\mathrm{Ker}}}
\def\Im{\mathop{\mathrm{Im}}}
\def\Coker{\mathop{\mathrm{Coker}}}
\def\Coim{\mathop{\mathrm{Coim}}}

% Boxtimes
%
\DeclareMathSymbol{\boxtimes}{\mathbin}{AMSa}{"02}

%
% Macros for moduli stacks/spaces
%
\def\QCohstack{\mathcal{QC}\!\mathit{oh}}
\def\Cohstack{\mathcal{C}\!\mathit{oh}}
\def\Spacesstack{\mathcal{S}\!\mathit{paces}}
\def\Quotfunctor{\mathrm{Quot}}
\def\Hilbfunctor{\mathrm{Hilb}}
\def\Curvesstack{\mathcal{C}\!\mathit{urves}}
\def\Polarizedstack{\mathcal{P}\!\mathit{olarized}}
\def\Complexesstack{\mathcal{C}\!\mathit{omplexes}}
% \Pic is the operator that assigns to X its picard group, usage \Pic(X)
% \Picardstack_{X/B} denotes the Picard stack of X over B
% \Picardfunctor_{X/B} denotes the Picard functor of X over B
\def\Pic{\mathop{\mathrm{Pic}}\nolimits}
\def\Picardstack{\mathcal{P}\!\mathit{ic}}
\def\Picardfunctor{\mathrm{Pic}}
\def\Deformationcategory{\mathcal{D}\!\mathit{ef}}
\usepackage{cleveref}
\crefname{originalremark}{original remark}{original remarks}
\Crefname{originalremark}{Original remark}{Original remarks}
\makeatletter
\newlabel{remark-direct-sum}{{0103}{source}{}{sourcepointer-1}{}}
\newlabel{remark-direct-sum@cref}{{[originalremark][0][]0103}{[1][0][]source}{}{}{}}
\makeatother
\crefname{originallemma}{original lemma}{original lemmas}
\Crefname{originallemma}{Original lemma}{Original lemmas}
\makeatletter
\newlabel{lemma-preadditive-zero}{{00ZZ}{source}{}{sourcepointer-2}{}}
\newlabel{lemma-preadditive-zero@cref}{{[originallemma][0][]00ZZ}{[1][0][]source}{}{}{}}
\makeatother
\title{An additive functor preserves direct sums and zero objects}
\author{The Stacks Project Authors}
\date{stacks-project:homology.tex:local-snapshot-20261012; GNU FDL 1.2 or later}
\begin{document}
\maketitle
\noindent\textit{Editorial note.} All delivered mathematical text is the author's own, copied from the pinned Stacks Project chapter file homology.tex; no statement or proof step was rewritten, reordered or removed.\par\medskip
\noindent\textit{Editorial note.} The proof is the author's complete argument for both assertions; every identity involved is either displayed in the counted span or named in the cited remark, and no step of the proof is left to the reader.\par\medskip
\noindent\textit{Editorial note.} Every labelled reference occurring inside the counted spans is carried in assembly.external\_references as a non-load-bearing navigation pointer (same-chapter ones at their definition site, cross-chapter ones at the invocation inside the counted span), so no load-bearing mathematics is replaced by a stub. No bibliography citation occurs inside the counted spans.\par\medskip
\noindent\textit{Editorial note.} The same-name check for this card scanned every staged Stacks chapter of this repository -- all sources/stacks-*/pinned-original-source.tex files, the card's own chapter excluded -- and the labels found are recorded in dedup.same\_name\_cross\_chapter on the card.\par\medskip
\noindent\textit{Editorial note.} The preamble is the chapter's own native preamble excerpt already used for this source: the comment-type reference, slogan and history environments (native preamble lines 8-13), the Xy-pic setup (native preamble lines 14-18) and the native theorem declarations with the mathematical macros (native preamble lines 154-217, the actual last line of the file).\par\medskip
\noindent\textit{Allowed standard prerequisites (editorial).}\begin{itemize}
\item Direct sums described by four morphisms. Same chapter, Remark remark-direct-sum. Supplies the list of identities satisfied by the two injections and the two projections of a direct sum, which is the data transported by the functor.
\item Characterizations of the zero object. Same chapter, Lemma lemma-preadditive-zero. Supplies the characterization of the zero object by which the second assertion is proved.
\end{itemize}\medskip
\ifdefined\hypertarget\hypertarget{sourcepointer-1}{}\fi
\noindent\textit{Source reference.} Original remark 0103: sources/stacks-homology/pinned-original-source.tex, source lines 372--372. Navigation pointer invoked at proof line 403; the counted span contains the pointer only and the cited result is not reproduced inside it, so the pointer is registered as a non-load-bearing navigation stub in the established form for Stacks material. No bibliography citation is involved.\par\medskip
\ifdefined\hypertarget\hypertarget{sourcepointer-2}{}\fi
\noindent\textit{Source reference.} Original lemma 00ZZ: sources/stacks-homology/pinned-original-source.tex, source lines 287--287. Navigation pointer invoked at proof line 412; the counted span contains the pointer only and the cited result is not reproduced inside it, so the pointer is registered as a non-load-bearing navigation stub in the established form for Stacks material. No bibliography citation is involved.\par\medskip
\begin{lemma}
\label{lemma-additive-additive}
Let $\mathcal{A}$, $\mathcal{B}$ be preadditive categories.
Let $F : \mathcal{A} \to \mathcal{B}$ be an additive functor.
Then $F$ transforms direct sums to direct sums and zero to zero.
\end{lemma}

\begin{proof}
Suppose $F$ is additive. A direct sum $z$
of $x$ and $y$ is characterized by having morphisms
$i : x \to z$, $j : y \to z$, $p : z \to x$ and
$q : z \to y$ such that $p \circ i = \text{id}_x$,
$q \circ j = \text{id}_y$, $p \circ j = 0$, $q \circ i = 0$
and $i \circ p + j \circ q = \text{id}_z$, according
to Remark \ref{remark-direct-sum}. Clearly $F(x), F(y), F(z)$
and the morphisms $F(i), F(j), F(p), F(q)$ satisfy exactly the
same relations (by additivity) and we see that $F(z)$ is
a direct sum of $F(x)$ and $F(y)$.
Hence, $F$ transforms direct sums to direct sums.

\medskip\noindent
To see that $F$ transforms zero to zero, use the
characterization (3) of the zero object in
Lemma \ref{lemma-preadditive-zero}.
\end{proof}

\end{document}
